\chapter{Estado del Arte}

Realizar la revisión del estado del arte con la metodología prisma, la cual vió en su espacio académico de investigación 1.
Resuma las investigaciones previas relacionadas con su tema.
Identifique las limitaciones y vacíos en la literatura actual que su investigación abordará.
Asegúrese de incluir las referencias más recientes y relevantes en su campo (ventana de 5 años)
Incluya una matriz de análisis de variables de comparación.
Nota: el uso de inteligencia artificial generativa esta permito siempre y cuando se mencione o referencie. Se debe evidenciar el mayor aporte al escrito realizado por usted. Tenga cuidado en verificar resúmenes y análisis extraídos por la IA, en ocasiones estos no dan cuenta de un análisis profundo.

\section{Categorías de Análisis}
\addcontentsline{toc}{section}{Introduction}
Describa cuantás y cuáles categorías de análsis escogio para el análisis de trabajos relacionados. 
\\

\subsection{Escriba el nombre de la primera CA}
Describa las categorías, cada una entre 4 a 10 párrafos, justificando con autores.

\subsection{Escriba el nombre de la segunda CA}
Describa las categorías, cada una entre 4 a 10 párrafos, justificando con autores.

\subsection{Escriba el nombre de la tercera CA}
Describa las categorías, cada una entre 4 a 10 párrafos, justificando con autores.

\section{Criterios}
Describa la importancia de los criterios en la revisión del estado del arte.
\subsection{Criterios de Inclusión}
xxxxxxx
\subsection{Criterios de Exclusión}
xxxxxxx

\section{Preguntas de investigación de la revisión sistemática}
Mencione las preguntas que estableció para analizar los artículos,

\section{Resultados}
1.Realicé análsis estadístico de la matriz general de artículos.
2.Deberá contestar cada una de las preguntas anteriores, referenciando autores y comparaciones.
3.Realicé un análisis descriptivo de la matríz secundaria y los RAE. Esto deberá hacerlo narrativo, en párrafos y referenciando.